\section{Knots, braids, and tangles in tori}

\begin{theorem}[A charge-coloured tangle invariant in $T^2\times I$]\label{thm:torus-tangle-charge-invariant}\notready\uses{thm:textile-link-on-torus, thm:overunder-negate}
\emph{(Open frontier --- stated, not formalized.)} The tangle-in-tori invariant is
a state-sum over the charge-coloured link of \cref{thm:textile-link-on-torus}.  Each
crossing resolution is a Belnap value, and the bracket is invariant under the
over/under involution $\sim$ from \cref{thm:overunder-negate}.  The residue is the
state-sum carrier on the torus link.
\end{theorem}

\begin{theorem}[Bobbin lace is a braid-word labelling of a torus graph]\label{thm:bobbin-lace-braid-word}\notready\uses{thm:torus-tangle-charge-invariant, thm:quasiperiodic-weave-recognizable}
\emph{(Open frontier --- stated, not formalized.)} A bobbin-lace ground is a
$2$-regular directed torus graph with vertices labelled by braid words.  In this
carrier, each vertex is a charge-permutation generator and admissibility is the
$\mathbb{Z}^2$ neighbourhood constraint.  The recognizable cover is supplied by
\cref{thm:quasiperiodic-weave-recognizable}; the residue is the braid-generator
action on the torus carrier.
\end{theorem}

\begin{theorem}[A bead-crochet rope is a torus knot]\label{thm:beadcrochet-rope-torus-knot}\notready\uses{thm:torus-tangle-charge-invariant, thm:mobius-band-holonomy}
\emph{(Open frontier --- stated, not formalized.)} A bead-crochet rope is a $1$D
thread walk embedded in the $\mathbb{Z}^2$ torus carrier by the helical offset
$\Phi(k)=(\lfloor k/m\rfloor,\,k+\alpha\lfloor k/m\rfloor\bmod m)$.  Closure
selects the $(p,q)$ torus-knot class.  A M\"obius rope adds the nontrivial
over/under holonomy of \cref{thm:mobius-band-holonomy}.  The residue is the
helical embedding $\mathbb{Z}\to\mathbb{Z}^2/\Lambda$ and its layer-group orbits.
\end{theorem}
